\section{Method}\label{sec:method}

\subsection{Problem Formalization}

We study test-time scaling for structured information extraction.
Given an input text $x$, a language model $\mathcal{M}$ generates $N$ independent candidate outputs $\{y_1, \ldots, y_N\}$ at temperature $T > 0$.
Each candidate $y_i$ is a structured object: an entity set $E_i = \{e_1, \ldots, e_{|E_i|}\}$ for named entity recognition (NER), or a relation triple set $R_i = \{(h, t, l)_1, \ldots\}$ for relation extraction (RE).
The central question is whether the redundancy across $N$ samples can be exploited to improve extraction quality.

We distinguish two objectives that prior work has conflated:
\begin{itemize}
    \item \textbf{Quality estimation (QE):} rank \emph{instances} by expected extraction quality.
    Given a set of instances $\{x^{(1)}, \ldots, x^{(M)}\}$, assign each a confidence score $c(x)$ such that $c(x)$ correlates with instance-level F1.
    This enables downstream applications such as quality-aware routing, selective prediction, and active verification.

    \item \textbf{Sample selection:} for a single instance $x$, rank the $N$ \emph{candidates} to select the best one.
    Assign each candidate a score $s(y_i \mid x)$ such that $\arg\max_i s(y_i \mid x)$ approximates the oracle choice $\arg\max_i \mathrm{F1}(y_i, y^*)$.
\end{itemize}

These two objectives impose different requirements on the confidence signal.
QE benefits from \textit{instance-level aggregation}: a signal that summarizes agreement patterns across all $N$ samples can distinguish easy instances (high agreement) from hard ones (low agreement), even though it assigns the same score to every candidate within an instance.
Selection, by contrast, requires \textit{per-sample discrimination}: the signal must assign \emph{different} scores to different candidates of the same instance.
A signal that aggregates across samples to produce a single instance-level score is structurally unable to perform selection, regardless of its correlation with quality.
This distinction is the key to understanding the correlation-selection gap we identify empirically (\S\ref{sec:experiments}).

\subsection{Confidence Signals}\label{sec:signals}

We evaluate five confidence signals drawn from three families.
Four signals aggregate across the $N$ samples to produce a single score per instance (instance-level); only log-probability provides per-sample scores.

\subsubsection{Structural Signals}

Structural signals operate directly on the extracted entity or relation sets, exploiting set-theoretic overlap.

\paragraph{Voting Confidence (VC).}
For each unique extraction $e$ in the union across all samples $U = \bigcup_{i=1}^{N} E_i$, we compute its frequency $\mathrm{freq}(e) = \sum_{i=1}^{N} \mathbf{1}[e \in E_i]$.
The voting confidence of an instance is the mean normalized frequency:
\begin{equation}
    \mathrm{VC}(x) = \frac{1}{|U|} \sum_{e \in U} \frac{\mathrm{freq}(e)}{N}
    \label{eq:vc}
\end{equation}
VC reflects how consistently the model extracts each item: high VC indicates that most extractions appear in most samples.
For RE, each extraction $e$ is a $(h, t, l)$ triple, and frequency is computed over exact triple matches.

\paragraph{Soft Jaccard (SJ).}
We measure pairwise set similarity across all sample pairs using the Jaccard index \cite{jaccard1901distribution}:
\begin{equation}
    \mathrm{SJ}(x) = \frac{2}{N(N-1)} \sum_{i < j} \frac{|E_i \cap E_j|}{|E_i \cup E_j|}
    \label{eq:sj}
\end{equation}
where intersection and union use exact span matching for NER and exact $(h, t, l)$ matching for RE.
When both $E_i$ and $E_j$ are empty, we define the Jaccard coefficient as 1.
SJ captures overall structural agreement: a score near 1 means the model produces nearly identical extraction sets across samples.

\subsubsection{Surface Signals}

Surface signals operate on the string representations of model outputs rather than on parsed structures.

\paragraph{Fleiss' $\kappa$ (FK).}
We treat the $N$ samples as independent annotators rating categories in a normalized text space.
Following \citet{fleiss1971agreement}, we compute:
\begin{equation}
    \mathrm{FK}(x) = \frac{\bar{P} - \bar{P}_e}{1 - \bar{P}_e}
    \label{eq:fk}
\end{equation}
where $\bar{P}$ is the mean observed agreement and $\bar{P}_e$ is the expected agreement under chance.
FK accounts for agreement expected by chance, providing a calibrated measure of inter-sample consensus.

\paragraph{Exact Match (EM).}
The fraction of sample pairs that are identical after normalization:
\begin{equation}
    \mathrm{EM}(x) = \frac{2}{N(N-1)} \sum_{i < j} \mathbf{1}\big[\mathrm{norm}(y_i) = \mathrm{norm}(y_j)\big]
    \label{eq:em}
\end{equation}
where $\mathrm{norm}(\cdot)$ applies whitespace and casing normalization.
EM is the coarsest signal: it is strictly binary per pair and collapses all partial agreement into disagreement.
As an instance-level aggregation, it is inherently unsuitable for sample selection.

\subsubsection{Model-Internal Signal}

\paragraph{Log-Probability (LP).}
The mean token-level log-likelihood under the model:
\begin{equation}
    \mathrm{LP}(y_i) = \frac{1}{|y_i|} \sum_{t=1}^{|y_i|} \log P_{\mathcal{M}}(y_{i,t} \mid y_{i,<t},\, x)
    \label{eq:lp}
\end{equation}
where $|y_i|$ is the number of tokens in $y_i$.
LP is the \textbf{only per-sample signal} in our comparison---it assigns a distinct score to each candidate $y_i$, making it the only signal structurally capable of sample selection.
However, constrained decoding (which forces valid output structure) compresses the token-level probability distribution, reducing score variance across candidates within an instance.
We quantify this compression via the \textit{tied fraction}: the proportion of instances where $\max_i \mathrm{LP}(y_i) - \min_i \mathrm{LP}(y_i) < \epsilon$ (with $\epsilon = 0.05$), indicating that LP cannot meaningfully discriminate among candidates.

\subsection{Evaluation Protocol}

\paragraph{Quality Estimation.}
We evaluate QE using two complementary metrics.
\textit{Spearman} $\rho$ measures rank correlation between each signal and instance-level F1 across all test instances.
\textit{AUROC} measures binary discrimination: whether the signal can distinguish instances above vs.\ below median F1.
Higher values indicate better quality estimation on both metrics.

\paragraph{Sample Selection.}
Selection quality is measured by \textbf{best-of-$N$ F1}: for each instance, the signal selects the candidate with the highest score, and we report the macro-averaged F1 of the selected candidates.
For the four instance-level signals (VC, SJ, FK, EM), which assign the same score to all $N$ candidates of an instance, selection reduces to random choice; we report the expected value over random selections.
The \textbf{oracle} upper bound selects the candidate with the highest F1, and the \textbf{greedy} baseline uses a single sample at $T{=}0$.

\paragraph{Gold-Filtered Evaluation.}
We exclude instances with no gold annotations from all evaluations, as these contribute undefined F1 scores and distort correlation estimates.
This ensures that our metrics reflect true signal quality rather than artifacts of empty-gold instances.
